\documentclass[conference]{IEEEtran}
\IEEEoverridecommandlockouts
\usepackage{cite}
\usepackage{amsmath,amssymb,amsfonts}
\usepackage{algorithmic}
\usepackage{graphicx}
\usepackage{textcomp}
\usepackage{xcolor}
\def\BibTeX{{\rm B\kern-.05em{\sc i\kern-.025em b}\kern-.08em
    T\kern-.1667em\lower.7ex\hbox{E}\kern-.125emX}}

\title{A Hybrid OCR Pipeline for Extracting Structured Information from Medical Prescriptions}

\author{
\IEEEauthorblockN{Medical Prescription OCR Research Team}
\IEEEauthorblockA{\textit{Medical Prescription OCR Project}\\
Educational Research Prototype}
\and
\IEEEauthorblockN{Sumit Mund}
\IEEEauthorblockA{\textit{Medical Prescription OCR Project}\\
Educational Research Prototype}
\and
\IEEEauthorblockN{Swapnaraj Mohanty}
\IEEEauthorblockA{\textit{Medical Prescription OCR Project}\\
Educational Research Prototype}
}

\begin{document}

\maketitle

\begin{abstract}
Medical prescriptions contain clinically important information such as medicine
names, dosage quantities, and administration frequencies, but their visual
formats vary considerably across printed and handwritten documents. This paper
presents the design of an educational medical-prescription optical character
recognition (OCR) pipeline that combines image preprocessing, EasyOCR, a
Tesseract fallback, and rule-based field extraction. The implementation converts
input images into multiple grayscale representations, performs OCR with a
GPU-aware EasyOCR reader, groups detected text boxes into lines, and invokes
Tesseract when the primary result is short or insufficiently confident. Extracted
text is then mapped to candidate medicine, dosage, and frequency fields using
regular expressions. Character error rate and word error rate are provided for
evaluation. The paper describes the implemented architecture, algorithms,
evaluation protocol, limitations, and safety requirements. No clinical
performance claims are made because the current repository does not contain a
completed, reproducible benchmark report.
\end{abstract}

\begin{IEEEkeywords}
medical prescription, optical character recognition, EasyOCR, Tesseract,
image preprocessing, information extraction
\end{IEEEkeywords}

\section{Introduction}
\label{sec:introduction}

Prescription digitization can reduce manual transcription effort and make
medication information easier to search and review. However, prescription
images present several challenges: variable illumination, skew, sensor noise,
different writing styles, mixed layouts, and ambiguous medical abbreviations.
Recognition errors are especially important because a visually plausible
output can still contain an incorrect medicine or dosage.

This project develops a modular OCR pipeline for prescription images. The
implemented system uses OpenCV for preprocessing, EasyOCR as the primary text
recognizer, Tesseract as a secondary recognizer, and lightweight regular
expressions for extracting structured candidates. The design prioritizes
inspectable intermediate stages and explicit confidence-based fallback behavior
over opaque end-to-end predictions. The implementation uses established
computer-vision and OCR components \cite{opencv,easyocr,tesseract}.

The contributions of this work are:
\begin{enumerate}
    \item a staged preprocessing pipeline that preserves intermediate image
    representations for inspection;
    \item a reusable OCR interface with lazy model loading and optional GPU
    execution;
    \item a fallback strategy that combines OCR confidence and recognized-text
    length;
    \item a transparent rule-based extractor for medicine, dosage, and
    frequency candidates; and
    \item an evaluation framework based on character and word error rates.
\end{enumerate}

The system is an educational prototype and is not intended for clinical
decision-making, dispensing, or autonomous prescription interpretation.

\section{Related Problem and Design Motivation}
\label{sec:motivation}

Prescription OCR is not equivalent to generic document transcription. The
system must preserve short tokens, units, abbreviations, and numbers while
handling text that may be handwritten or spatially irregular. A single OCR
engine can fail on a particular image condition, so the project uses two
recognizers with different operational characteristics. EasyOCR provides
bounding boxes and confidence values suitable for line reconstruction, while
Tesseract supplies an independent fallback path.

The implementation also separates recognition from field extraction. This
separation makes it possible to inspect raw OCR text before applying
domain-specific heuristics, and it allows future work to replace the extractor
without changing image processing or OCR interfaces.

\section{System Architecture}
\label{sec:architecture}

The complete workflow is shown conceptually in Fig.~\ref{fig:pipeline}.
Images are indexed and opened as RGB data, converted to grayscale, optionally
deskewed, denoised, and binarized. OCR is then attempted using EasyOCR.
Tesseract is invoked only when the primary result does not meet the configured
confidence and length requirements. Finally, the selected text is passed to the
rule-based field extractor.

\begin{figure}[t]
\centering
\fbox{\parbox{0.92\columnwidth}{
\centering
\vspace{0.35cm}
Image loading $\rightarrow$ preprocessing $\rightarrow$ EasyOCR\\
$\downarrow$ low confidence/short text $\rightarrow$ Tesseract fallback\\
$\rightarrow$ field extraction $\rightarrow$ CER/WER evaluation
\vspace{0.35cm}
}}
\caption{Implemented prescription OCR workflow.}
\label{fig:pipeline}
\end{figure}

The repository is organized into source modules for data access, preprocessing,
OCR, extraction, and metrics. Three notebooks support data preparation,
preprocessing inspection, and OCR evaluation. Gold-label CSV files are stored
under the \texttt{data} directory.

\section{Methodology}
\label{sec:methodology}

\subsection{Data Access and Image Representation}
\label{sec:data}

The data utility module recursively indexes images with PNG, JPEG, and WebP
extensions. It maps each base filename to a path and opens images using
Pillow, converting them to RGB. Gold annotations are represented as grouped
records containing the fields \texttt{medicine}, \texttt{dosage}, and
\texttt{frequency}.

The current project includes development and test gold-label CSV files and
notebooks for generating dataset splits. The repository documentation describes
a synthetic prescription dataset workflow, but this paper does not treat the
described split sizes as measured experimental results.

The supplied training-validation notebook validates the EasyOCR trainer
manifests before model training. It checks that each label file has the
\texttt{filename} and \texttt{words} columns, that the number of labels equals
the number of PNG images, and that every referenced image exists. The validated
manifest counts are shown in Table~\ref{tab:dataset}.

\begin{table}[t]
\centering
\caption{Validated training and validation manifests.}
\label{tab:dataset}
\begin{tabular}{lrr}
\hline
\textbf{Split} & \textbf{Samples} & \textbf{Validation checks}\\
\hline
\texttt{medical\_train} & 3,120 & 3\\
\texttt{medical\_val} & 780 & 3\\
\hline
\end{tabular}
\end{table}

These counts describe manifest validation only; they are not OCR accuracy,
loss, or generalization results.

\subsection{Image Preprocessing}
\label{sec:preprocessing}

The preprocessing function returns a dictionary of stages so that downstream
experiments can compare OCR performance on different representations. A
Pillow image is converted to a BGR NumPy array before OpenCV operations.
The image is converted to grayscale and resized so that its longest side is
approximately 1600 pixels while preserving its aspect ratio. Cubic
interpolation is used when enlarging and area interpolation when reducing.
Images whose scale is already within five percent of the target are left
unchanged.

The resized grayscale image is optionally deskewed. Deskewing estimates the
orientation of foreground pixels using a minimum-area rectangle and ignores
tiny or suspicious rotations. Non-Local Means denoising then reduces speckles
and sensor noise. Finally, Gaussian adaptive thresholding produces a binary
image designed to handle uneven lighting and shadows. The default stage keys
are \texttt{gray}, \texttt{resized}, \texttt{denoised}, and \texttt{binary};
\texttt{deskewed} is included when deskewing is enabled.

\subsection{EasyOCR Recognition}
\label{sec:easyocr}

EasyOCR is loaded lazily through a module-level reader cache. On the first
request, the implementation detects CUDA availability through PyTorch unless
the caller explicitly supplies a GPU choice. Subsequent requests reuse the
reader, avoiding repeated model initialization.

EasyOCR returns detected regions as bounding boxes, text, and confidence
values. The intended line-grouping procedure computes the horizontal position,
vertical center, and height of each region. Regions with sufficiently similar
vertical centers are grouped into a line, and each line is sorted from left to
right. The grouped lines are joined with newline separators. The EasyOCR
confidence reported by the pipeline is the arithmetic mean of the individual
region confidences.

\subsection{Tesseract Fallback}
\label{sec:tesseract}

Tesseract is configured with:

\begin{verbatim}
--oem 3 --psm 6
\end{verbatim}

The Tesseract path obtains word confidence values through
\texttt{image\_to\_data} and obtains line text through
\texttt{image\_to\_string}. Tesseract confidence percentages are normalized to
the interval $[0,1]$ by division by 100. Empty or invalid confidence entries
are excluded.

The fallback controller first executes EasyOCR. It returns that result when
confidence is at least 0.5 and the recognized text contains at least 20
characters by default. Otherwise, it executes Tesseract and chooses the result
with the larger score:

\begin{equation}
S(r) = C(r)\min(L(r), 200),
\label{eq:fallback}
\end{equation}

where $C(r)$ is the normalized confidence and $L(r)$ is the length of the
recognized text. This score rewards confidence while preventing unusually long
outputs from dominating without bound. When fallback processing is used, the
reported elapsed time includes both OCR calls and the result includes a
\texttt{fallback\_used} flag.

\subsection{Rule-Based Information Extraction}
\label{sec:extraction}

The extractor removes blank lines and examines each remaining line. Dosage
candidates are identified with a case-insensitive pattern consisting of a
number followed by one of the units \texttt{mg}, \texttt{ml}, \texttt{mcg},
\texttt{g}, or \texttt{iu}. Frequency candidates include patterns such as
\texttt{2x daily}, \texttt{daily}, \texttt{times}, \texttt{bd}, \texttt{tds},
\texttt{od}, and \texttt{hs}. Lines longer than three characters that contain no
digits are treated as medicine-name candidates.

The extractor returns three lists:

\begin{verbatim}
{
    "medicine": [...],
    "dosage": [...],
    "frequency": [...]
}
\end{verbatim}

These are candidates rather than validated medication instructions. The
current implementation does not link a dosage and frequency to a specific
medicine, correct OCR spelling, or verify a medicine against a clinical
database.

\section{Evaluation Protocol}
\label{sec:evaluation}

\subsection{Recognition Metrics}

The project provides character error rate (CER) and word error rate (WER).
CER is computed from normalized Levenshtein distance \cite{levenshtein}:

\begin{equation}
\mathrm{CER}(p,t)=\frac{d_{\mathrm{Lev}}(p,t)}{|t|},
\label{eq:cer}
\end{equation}

where $p$ is the prediction and $t$ is the reference text. WER is estimated
by splitting both strings into words and computing one minus the
\texttt{SequenceMatcher} similarity ratio:

\begin{equation}
\mathrm{WER}(p,t)=1-\mathrm{SequenceMatcher}(p,t).
\label{eq:wer}
\end{equation}

Lower values indicate better agreement. Empty-reference behavior is explicitly
defined in the implementation: an empty prediction receives zero error, while
a non-empty prediction receives one.

\subsection{Proposed Experimental Comparisons}

The notebooks support comparisons among:

\begin{enumerate}
    \item EasyOCR on each preprocessing stage;
    \item Tesseract on each preprocessing stage;
    \item EasyOCR alone versus confidence-triggered fallback;
    \item CPU execution versus GPU execution when available; and
    \item raw OCR output versus extracted field candidates.
\end{enumerate}

For each image, the evaluation should record the selected engine, confidence,
elapsed time, OCR text, CER, WER, and whether fallback was used. Results should
be aggregated over a fixed development set and confirmed on a held-out test
set. Because the current repository documentation does not provide a completed
metric table, numerical performance values are intentionally not reported in
this paper.

\section{Implementation Status}
\label{sec:status}

The data, preprocessing, extraction, and metric modules describe a coherent
prototype. However, the current \texttt{src/ocr.py} contains an indentation
error in the \texttt{group\_lines} function and cannot currently be compiled.
The file also contains an incomplete line-grouping body with references to
\texttt{bbox} and \texttt{text} before the required result iteration. Therefore,
the architecture and algorithm described in Sections~\ref{sec:easyocr} and
\ref{sec:tesseract} represent the intended implementation, but end-to-end OCR
execution requires correcting this source defect first.

\noindent This status is important for reproducibility: no claim in this paper should be
interpreted as a successful benchmark until the OCR module is repaired and the
evaluation notebooks are executed on a documented environment.

\section{Limitations and Safety}
\label{sec:limitations}

The prototype has several technical limitations:

\begin{itemize}
    \item Mean OCR confidence is not a guarantee of correct medication text.
    \item The fallback score is heuristic and does not measure semantic
    correctness.
    \item Regular expressions may miss uncommon units, abbreviations, and
    prescription layouts.
    \item Medicine candidates are not validated or linked to dosage and
    frequency fields.
    \item Base-filename indexing can overwrite files with duplicate names.
    \item Handwriting and low-quality images may require specialized models or
    human review.
    \item The repository does not currently contain a complete readable test
    suite or a reproducible final benchmark.
\end{itemize}

This system must not be used to prescribe, dispense, or change medication.
Every recognized medicine, dosage, and frequency must be verified against the
original prescription by a qualified reviewer. A production system would also
require privacy controls, audit logs, calibrated confidence estimates,
clinical validation, error analysis by prescription type, and explicit
human-in-the-loop approval.

\section{Future Work}
\label{sec:future}

Future development should first repair and test the OCR line-grouping logic.
Additional work should include:

\begin{enumerate}
    \item unit tests for every preprocessing and extraction boundary case;
    \item deterministic dataset manifests and duplicate-filename detection;
    \item field-level precision, recall, and F1 evaluation;
    \item confidence calibration and threshold tuning on a development set;
    \item vocabulary-assisted medicine correction using the available medicine
    lists;
    \item association of each medicine with its dosage and frequency;
    \item specialized handwriting recognition and layout analysis;
    \item privacy-preserving storage and deletion policies; and
    \item prospective human review studies before any clinical deployment.
\end{enumerate}

\section{Conclusion}
\label{sec:conclusion}

This paper presented an inspectable hybrid OCR architecture for medical
prescription images. The implementation combines configurable OpenCV
preprocessing, GPU-aware EasyOCR, Tesseract fallback, simple structured-field
extraction, and CER/WER evaluation. Its modular design supports controlled
experiments and clear diagnosis of recognition failures. At its current stage,
the system should be treated as an educational prototype rather than a
clinical tool. Correcting the OCR module, adding reproducible tests, and
reporting measured field-level results are necessary steps before the approach
can be evaluated for practical use.

\begin{thebibliography}{00}

\bibitem{opencv}
G. Bradski, ``The OpenCV Library,'' Dr. Dobb's Journal of Software Tools,
2000.

\bibitem{easyocr}
JaidedAI, ``EasyOCR: Ready-to-use OCR with 80+ supported languages,''
GitHub repository, accessed Nov. 5, 2026.

\bibitem{tesseract}
R. Smith, ``An Overview of the Tesseract OCR Engine,'' in
\emph{Proc. Ninth International Conference on Document Analysis and
Recognition}, 2007, pp. 629--633.

\bibitem{levenshtein}
V. I. Levenshtein, ``Binary codes capable of correcting deletions, insertions,
and reversals,'' \emph{Soviet Physics Doklady}, vol. 10, no. 8, pp. 707--710,
1966.

\end{thebibliography}

\end{document}
